\begin{figure}[h]
\sporkSemanticsStuff
\hrule
\begin{mathpar}
\inferrule*[right=\sporkERead]{
  \sporkvalue = \sporksigma(\sporkx, \sporkz^*)
}
{
  \sporkRule{\sporkSub{\sporkx[\sporkz^*]}{\sporkvalue}, \sporksigma, \sporktr^*}
}
\inferrule*[right=E-Negate]{
  \sporkRule{\sporkSub{\sporkexpr_1}{\sporkexpr_2}, \sporksigma, \sporkSub{\sporktr^*_1}{\sporktr^*_2}}
}
{
  \sporkRule{-\sporkSub{\sporkexpr_1}{\sporkexpr_2}, \sporksigma, \sporkSub{\sporktr^*_1}{\sporktr^*_2}}
}\\
\inferrule*[right=E-BinOpLhs]{
  \oplus \in \sporkBinOp \\\\
  \sporkRule{\sporkSub{\sporkexpr_1}{\sporkexpr_2}, \sporksigma, \sporkSub{\sporktr^*_1}{\sporktr^*_2}}
}
{
  \sporkRule{\sporkSub{\sporkexpr_1}{\sporkexpr_2} \oplus \sporkexpr, \sporksigma, \sporkSub{\sporktr^*_1}{\sporktr^*_2}}
}
\inferrule*[right=E-BinOpRhs]{
  \oplus \in \sporkBinOp \\\\
  \sporkRule{\sporkSub{\sporkexpr_1}{\sporkexpr_2}, \sporksigma, \sporkSub{\sporktr^*_1}{\sporktr^*_2}}
}
{
  \sporkRule{\sporkvalue \oplus \sporkSub{\sporkexpr_1}{\sporkexpr_2}, \sporksigma, \sporkSub{\sporktr^*_1}{\sporktr^*_2}}
}\\
\inferrule*[right=E-ExternArg]{
  \sporkRule{\sporkSub{\sporkexpr_1}{\sporkexpr_2}, \sporksigma, \sporkSub{\sporktr^*_1}{\sporktr^*_2}}
}
{
  \sporkRule{\sporkextern(\sporkvalue^*, \sporkSub{\sporkexpr_1}{\sporkexpr_2}, \sporkexpr^*), \sporksigma, \sporkSub{\sporktr^*_1}{\sporktr^*_2}}
}
\end{mathpar}
\caption{Expression Semantics.
``Read only'' expressions must be carefully stepped to preserve trace effects.
We assume constant arithmetic/logic, including externs,
e.g. axioms \mbox{$\sporkRule{\sporkSub{2 \texttt{+} 2}{4}, \sporksigma, \sporktr^*}$},
\mbox{$\sporkRule{\sporkSub{1 \texttt{==} 1}{\sporkTrue}, \sporksigma, \sporktr^*}$},
and \mbox{$\sporkRule{\sporkSub{\texttt{exp2}(3)}{8}, \sporksigma, \sporktr^*}$}.
Ill-typed combinations (e.g.\ \texttt{True + 1}) have no axiom, so they are stuck.
Note \sporkERead\ cannot operate on $\#\sporkwinAnn$ annotated windows, which are first stepped by \sporkTrRead\ (Figure~\ref{fig:SyncTraceSemantics}).
}
\label{fig:ExprSemantics}
\hrule
\end{figure}
